%!TEX ROOT = thesis.tex

% =========================================================
% CHAPTER 1: INTRODUCTION
% =========================================================
\chapter{INTRODUCTION}
\label{chap:intro}

\section{Background of the Study}

\subsection{Clinical Significance of Cardiovascular Diseases}
Cardiovascular diseases (CVDs) remain one of the most important global causes of mortality. According to the World Health Organization (WHO), CVDs caused an estimated 19.8 million deaths in 2022, corresponding to approximately 32\% of all global deaths \citep{who2025cvd}. This figure is not used in this thesis as a direct claim about the performance of any computational model. Rather, it establishes the clinical context in which efficient and reproducible cardiac image analysis tools are valuable. Given this burden, even a modest improvement has practical relevance if it makes cardiac assessment more dependable or easier to deliver at scale.

Cardiac Magnetic Resonance (CMR) imaging offers a non-invasive means of examining cardiac anatomy and function \citep{bernard2018deep}. It allows clinicians to obtain functional indices including Ejection Fraction (EF), End-Diastolic Volume (EDV), End-Systolic Volume (ESV), and Myocardial Mass. Reliable calculation of these indices requires separate masks of the myocardium (MYO) and the two ventricular cavities, LV and RV. Automated delineation consequently forms a core step in computer-assisted cardiac assessment. The experiments in this project first assess spatial segmentation at ED and ES. The selected Nano-Mamba checkpoint is subsequently run over every phase of each complete cine sequence to obtain phase timing, EF, ventricular volume trajectories, and global segmentation-derived motion surrogates.

The connection between a segmentation mask and a quantitative cardiac index can be illustrated with ventricular volume. If $V_{ED}$ is the ventricular volume at end-diastole and $V_{ES}$ is the volume at end-systole, ejection fraction is given by:
\begin{equation}
    EF = \frac{V_{ED}-V_{ES}}{V_{ED}}\times 100\%.
\end{equation}
The formula makes the consequence of a segmentation error explicit. A small boundary error can change the estimated cavity volume, and the resulting volume error can propagate into EF. In the completed full-cine analysis, the same relationship is evaluated at the manually labelled ED/ES endpoints and used to construct a complete segmentation-derived LV-volume trajectory.

\subsection{The Bottleneck of Manual Delineation and Traditional Methods}
Manual delineation of cardiac structures requires expert knowledge and substantial effort. A cine CMR study may contain many slices and cardiac phases. The ACDC challenge paper identifies both the diagnostic importance of accurately segmenting RV, MYO, and LV and the possibility of inter-observer variation in manual annotations \citep{bernard2018deep}. No single annotation-time estimate is assumed here because the available evidence does not support one. The relevant limitation is operational: a trained reader must contour each case, and the workload grows with the number of studies.

Manual contouring also involves judgement rather than simple tracing. Papillary muscles can obscure the boundary of the ring-shaped myocardium, and the RV has both a thin wall and marked geometric variation. Basal and apical slices introduce another ambiguity because cardiac structures appear or disappear gradually as the imaging plane changes. Two experts may therefore draw different contours on the same image. Automation does not replace that expertise; it provides a consistent initial result that a clinician can inspect before measurements are made.

Earlier segmentation pipelines used methods such as thresholding, active contours, level sets, and graph cuts. Their behaviour often depends on a chosen feature set, an initial contour, or a particular model of image intensity. CMR does not satisfy these conditions reliably: scanner-dependent contrast, partial-volume effects, papillary muscles, and indistinct boundaries all affect the image \citep{chen2020cardiacreview}. These limitations helped shift the field towards models whose representations are fitted from annotated data.

\subsection{Deep Learning for Cardiac Segmentation}
Deep learning changed this workflow by learning image features during training. Fully Convolutional Networks (FCNs) showed that a network could produce a dense semantic label map \citep{long2015fully}. U-Net then paired an encoder and decoder with skip connections, a design that became widely used in biomedical segmentation \citep{ronneberger2015u}; 3D U-Net extended it to volumetric convolutions \citep{cicek20163d}. The present work builds on these established architectures and does not claim them as new contributions.

A convolutional kernel sees only a local neighbourhood at each operation. For example, a $3\times3\times3$ kernel initially combines information from nearby voxels; a wider effective field emerges only after several layers are stacked. Cardiac MRI nevertheless requires both fine boundary evidence and knowledge of how the chambers and myocardium are arranged. The network must therefore separate myocardium from blood pool along local boundaries without losing the relative spatial organization of RV, MYO, and LV.

Transformer-based architectures address long-range interactions through self-attention \citep{vaswani2017attention,dosovitskiy2020image}. Medical segmentation models such as UNETR and Swin UNETR demonstrate the usefulness of global or window-based token interactions in volumetric segmentation tasks \citep{hatamizadeh2022unetr,hatamizadeh2022swin}. However, self-attention can increase memory and computational requirements, which motivates research into more efficient sequence modeling alternatives. This thesis does not argue that Transformers are unsuitable for medical imaging. Instead, it investigates whether a lighter sequence-inspired bottleneck can be useful when the hardware environment is constrained.

\subsection{Motivation for a Mamba-Inspired Bottleneck}
Selective State Space Models (SSMs), including Mamba, have recently attracted attention because they can model long sequences with linear-time complexity \citep{gu2023mamba}. This thesis investigates whether a lightweight Mamba-inspired sequence bottleneck can be integrated into a 3D U-Net-style architecture for cardiac MRI segmentation. The implemented model serializes compressed 3D volumetric features into a one-dimensional token sequence, applies a lightweight gated sequence module inspired by Mamba, and reshapes the result back into volumetric form.

The motivation is based on the accuracy-efficiency trade-off. A model with the highest Dice score is not necessarily the most suitable model for all environments if it requires more memory, more parameters, or longer inference time. Conversely, a very small model is not useful if it fails to segment clinically relevant structures. This thesis therefore evaluates Nano-Mamba U-Net not only by segmentation accuracy but also by parameter count and inference speed. The research question is whether a compact sequence-inspired bottleneck can produce competitive validation performance while maintaining a small model size.

The completed empirical design has two linked stages. In the historical six-model experiment, labelled ED and ES volumes are processed as independent 3D cases and tensor depth contains anatomical slices rather than time. In a separate post-hoc stage, the selected spatial Nano-Mamba checkpoint is applied to every frame of each validation cine. Fixed adjacent-frame probability fusion is then evaluated against frame-wise inference. From the resulting native-grid masks, the analysis derives ventricular volume trajectories together with ED/ES timing and EF. Motion is summarized separately through a global myocardial radial-distance surrogate and LV-centroid displacement. The project therefore completes a segmentation-derived spatio-temporal analysis, while stopping short of claiming a learned temporal Mamba network, optical flow, dense deformation, or regional strain.

\section{Scope of the Study}
The scope of this research is defined as follows:
\begin{itemize}
    \item \textbf{Dataset Boundary:} The empirical evaluation uses the Automated Cardiac Diagnosis Challenge (ACDC) dataset \citep{bernard2018deep}. The segmentation targets are RV, MYO, and LV.
    \item \textbf{Task Boundary:} The trained backbone performs four-class 3D cardiac MRI segmentation and treats ED/ES cases independently. A separate full-cine workflow evaluates frame-wise inference and fixed temporal probability fusion across 550 validation frames, then derives global volume, phase, EF, centroid-displacement, and myocardial-radius trajectories. Slice depth is never treated as time.
    \item \textbf{Implementation Boundary:} The proposed Nano-Mamba U-Net uses a Mamba-inspired bottleneck implemented in PyTorch. It should not be interpreted as a full reproduction of the original hardware-aware selective scan implementation.
    \item \textbf{Evaluation Boundary:} The main reported results are based on a deterministic patient-level 80/20 validation split of the ACDC training cohort. The study does not claim performance on the hidden official ACDC test set.
    \item \textbf{Hardware Boundary:} Training, validation, and computational profiling were conducted on a single consumer-grade NVIDIA RTX 4060 GPU with 8GB VRAM. Reported FPS is specific to the implemented batch-one benchmark and is not treated as a universal property of an architecture.
\end{itemize}
The scope statements above define how the results are to be read. This is a reproducible project-level investigation; it is neither a claim of clinical readiness nor evidence that the proposed model outperforms every cardiac segmentation method.

\section{Organization of the Thesis}
The thesis proceeds from context to evidence. Chapter 1 introduces the clinical and computational motivation, Chapter 2 defines the research problem, and Chapter 3 states the objectives and questions. The relevant convolutional, attention, residual, and state-space literature is reviewed in Chapter 4. Chapter 5 then records the data, preprocessing, model architecture, training, and evaluation protocol. Chapters 6 and 7 present and interpret the validation results, respectively, before Chapter 8 draws the conclusions and outlines future work.

% =========================================================
% CHAPTER 2: RESEARCH PROBLEM
% =========================================================
\chapter{RESEARCH PROBLEM}
\label{chap:problem}

An accurate segmentation score is necessary, but it is not the only design consideration for automated cardiac MRI. The model also has to run within available computing resources and handle the anatomical variation present across patients. This study therefore examines three connected difficulties in building a lightweight volumetric segmenter: limited contextual reach, computational cost, and variable cardiac anatomy.

\section{Local Context Modeling in Lightweight 3D CNNs}
3D CNNs process neighbouring voxels together, which is useful for volumetric segmentation. Their reach at one layer is nevertheless limited: a standard $3 \times 3 \times 3$ kernel only observes the surrounding voxels. Deeper stacks expand that reach, but depth and channel width are precisely what a small model must restrict. As a result, relationships across the heart may be represented incompletely, especially on basal and apical slices where the RV, MYO, and LV contours are uncertain. At such locations, similar local intensities can belong to different structures, so anatomical position and the surrounding chamber arrangement provide useful evidence beyond the immediate kernel.

Convolution remains the main operation in the proposed U-Net and is not being rejected. The practical issue is how much context survives when a small network compresses its features. Here, the bottleneck grid is written as a raster sequence and passed through a compact gated transformation. Mixing occurs over a kernel of three positions; the implemented block has neither selective state propagation nor direct all-token interaction. Its deliberately narrow role is why it supplements the compressed bottleneck instead of replacing the full U-Net with a large attention architecture. Operating after compression also keeps the sequence short enough for the available hardware. The experiment consequently evaluates a small addition to the convolutional backbone rather than a competing end-to-end sequence model.

\section{Computational Cost of Larger Segmentation Models}
Additional layers, channels, residual blocks, or attention operations can give a segmentation model more capacity. They also tend to increase its parameter count and may raise memory use or inference cost. Where hardware is limited, a small loss in Dice can be acceptable if it buys a substantial reduction in model size. The comparison in this thesis is therefore concerned with the observed accuracy--efficiency balance, not with Dice in isolation. Neither side of that balance is assumed in advance: segmentation quality and resource indicators are reported together for the executed models.

For that reason, the evaluation places parameter count and measured inference speed beside Dice. The term \emph{lightweight} is used here in a specific sense: the model has the lowest reported trainable-parameter count and competitive batch-one throughput on the recorded hardware. It is not claimed to lead simultaneously in memory, speed, and accuracy. Thus, the 1.46M-parameter Nano model can still be useful when a larger model attains a higher Dice, provided that the saving is material. All baselines and ablations are compared under the same data split and training protocol so that this trade-off can be examined consistently.

\section{Pathological and Anatomical Variability}
Besides normal cases, ACDC contains myocardial infarction, hypertrophic cardiomyopathy, dilated cardiomyopathy, and right-ventricular abnormality \citep{bernard2018deep}. These diagnoses introduce differences in chamber size, myocardial thickness, and boundary appearance. The validation data therefore include both normal hearts and anatomy altered by disease.

Anatomical variation is especially troublesome for the RV and MYO masks. The RV is thin-walled and irregular, while the myocardium occupies a narrow ring. A displacement of only a few voxels can therefore affect their overlap scores more strongly than it would affect a larger cavity. Reporting RV, MYO, and LV separately makes this behaviour visible instead of hiding it inside one mean Dice value. It also shows whether a change in the overall average reflects broad improvement or is driven mainly by the easier ventricular cavity.

\section{From Endpoint Segmentation to Cardiac-Cycle Analysis}
Manual labels at ED and ES make endpoint segmentation evaluation possible, but they do not show what the model predicts between those phases. Applying the spatial checkpoint to every cine frame fills that gap with a sequence of masks. Because the frames were not trained jointly, however, adjacent predictions can fluctuate. Examining the complete curve makes those jumps visible and permits phase and functional measurements that cannot be obtained from one isolated endpoint mask. The research problem is thus to connect the validated spatial checkpoint to a transparent time-series analysis while keeping it distinct from a learned four-dimensional network. The approach used here applies fixed circular probability fusion and assesses endpoint overlap, EF, phase timing, curve smoothness, and global motion surrogates. Manual endpoint labels provide partial validation, while the absence of intermediate-frame labels remains a stated limitation.

\section{Research Gap}
Existing segmentation backbones such as 3D U-Net, Attention U-Net, and SegResNet provide strong baselines, but they represent different points on the accuracy-efficiency spectrum. The first research gap addressed in this thesis is whether a compact Mamba-inspired bottleneck can provide competitive validation performance with fewer parameters than conventional or residual alternatives, while remaining feasible on a consumer GPU. The second is whether the selected spatial checkpoint can support a reproducible full-cine, segmentation-derived function and global motion analysis, and whether a fixed temporal fusion produces a measurable change relative to independent frame-wise inference.

The gap is not a claim that previous methods are inadequate in all settings. SegResNet, for example, remains a strong residual convolutional baseline. The gap is narrower and more appropriate for this project: there is value in testing a compact sequence-inspired bottleneck under a rigorous patient-level validation protocol, especially when earlier exploratory scripts may overestimate performance by evaluating on non-held-out data. This thesis therefore emphasizes reproducibility and honest comparison as much as architectural novelty.

% =========================================================
% CHAPTER 3: RESEARCH OBJECTIVES
% =========================================================
\chapter{RESEARCH OBJECTIVES}
\label{chap:objectives}

The aim of this research is to design, implement, and evaluate a lightweight Mamba-inspired U-Net architecture for cardiac MRI segmentation and to determine how its segmentations can support a bounded full-cine cardiac motion and function analysis. Empirical validation is central to the study: parameter efficiency is assessed alongside established baselines, and spatial learning is kept conceptually separate from the post-hoc temporal stage. The objectives connect directly to the experiments reported in Chapter 6.

\section{Specific Objectives}
\begin{enumerate}
    \item To implement a Nano-Mamba U-Net architecture that inserts a lightweight Mamba-inspired sequence bottleneck into a 3D U-Net-style encoder-decoder network.
    \item To evaluate the segmentation performance of the proposed model on RV, MYO, and LV using Dice Similarity Coefficient on a deterministic patient-level validation split of the ACDC training cohort.
    \item To compare the proposed model with 3D U-Net, Attention U-Net, SegResNet, and ablation variants using the same split, preprocessing, epoch budget, optimizer family, and validation protocol, with model-specific batch sizes disclosed as a limitation.
    \item To quantify computational efficiency using reported trainable-parameter count and batch-one inference speed.
    \item To apply the selected Nano-Mamba checkpoint to complete validation cine sequences, derive segmentation-based cardiac volume, phase, EF, and global motion trajectories, and quantify what fixed adjacent-frame probability fusion changes relative to independent frame-wise inference.
\end{enumerate}

The first objective is implemented through the model architecture described in Chapter 5. The second and third objectives are addressed through the six-model quantitative validation table in Chapter 6. The fourth objective is addressed by reporting trainable parameters and inference speed together with Dice. The fifth objective is addressed by the independently audited 20-patient, 550-frame full-cine analysis. This mapping ensures that every stated objective is answered by an implemented experiment.

\section{Research Questions}
\begin{enumerate}
    \item How can compressed 3D cardiac MRI feature maps be serialized and processed by a Mamba-inspired bottleneck within a U-Net-style segmentation architecture?
    \item Does the proposed Nano-Mamba U-Net achieve competitive held-out validation Dice compared with 3D U-Net, Attention U-Net, SegResNet, and ablation variants?
    \item What accuracy-efficiency trade-off is observed when comparing Nano-Mamba U-Net with heavier or non-Mamba alternatives in terms of Dice score, parameter count, and inference speed?
    \item How can the selected spatial segmentation checkpoint support full-cine cardiac function and global motion analysis, and what measurable effect does fixed adjacent-frame probability fusion have relative to frame-wise inference?
\end{enumerate}

These research questions are intentionally phrased in a measurable way. The first question is answered by architecture implementation rather than by a numerical score alone. The second question is answered by held-out validation Dice. The third question is answered by jointly interpreting Dice, parameter count, and speed. The fourth is answered by complete cine coverage, endpoint validation, functional and motion-derived quantities, and paired temporal-fusion comparisons. Each answer must therefore follow from an implemented analysis; neither absolute superiority nor dense motion tracking can be inferred from these questions.
